\documentclass{ctexart}
\usepackage{geometry}
\usepackage{fancyhdr}
\usepackage{graphicx}
\usepackage{booktabs}
\usepackage{amsmath}
\usepackage{tikz}
\usepackage{array}
\xeCJKsetup{CJKmath=true} 
\usepackage{zhnumber} % change section number to chinese
\renewcommand\thesection{\zhnum{section}}
\renewcommand \thesubsection {\arabic{subsection}}
\CTEXsetup[format={\Large\bfseries}]{section}

\geometry{
    a4paper,
    left=3.18cm,
    right=3.18cm,
    top=3.04cm,
    bottom=3.04cm
}

\pagestyle{fancy}
\fancyhf{}
\renewcommand{\headrulewidth}{0.7pt} % 设置页眉横线粗细
\fancyhead[L]{\kaishu\large 大学物理实验报告} % 在左侧设置页眉文字
\fancyhead[R]{\kaishu\large 哈尔滨工业大学(深圳) } % 在右侧设置页眉文字
\fancyfoot[R]{\thepage} % 将页数放在右下角


\setlength\headwidth{\textwidth}

\begin{document}

\noindent
\begin{center}
\textbf{
\begin{tabular}{p{2.4cm}p{2.4cm}p{4cm}p{4cm}}
    班级 \hrulefill & 学号 \hrulefill & 姓名 \hrulefill & 教师签字 \hrulefill \\
\end{tabular}
\begin{tabular}{p{6cm}p{3.6cm}p{3.6cm}}
    实验日期 \hrulefill & 预习成绩 \hrulefill & 总成绩 \hrulefill
\end{tabular}
{\noindent}	 \rule[-10pt]{\textwidth}{0.7pt}
}\end{center}

\begin{center}
    \Large \textbf{实验内容 \underline{用惠斯通电桥测电阻}}
\end{center}



\section{实验目的}

\vspace{6\baselineskip}

\section{实验预习}

绘制惠斯通电桥的电路图，并说明平衡时满足的条件。

\newpage
\section{实验现象及数据记录}

\subsection{惠斯通电桥测量电阻}

\begin{table}[!h]
    \centering
    \renewcommand{\arraystretch}{1.5} % 表格行高倍数
    \begin{tabular}{|c|m{1.3cm}<{\centering}|m{1.3cm}<{\centering}|m{1.3cm}<{\centering}|m{1.3cm}<{\centering}|m{1.3cm}<{\centering}|m{1.3cm}<{\centering}|}
        \hline
        电阻（阻值） & $N$ & $R_s (\varOmega)$ & $R_x (\varOmega)$ & $\Delta R_s (\varOmega)$ & $\Delta n$（格） & $S$（格） \\
        \hline
        1 $ k\varOmega $ & 1 & & & & & \\
        \hline
        100 $ \varOmega $ & 1 & & & & & \\
        \hline
    \end{tabular}
\end{table}

\subsection{惠斯通电桥灵敏度测量(1 $ k\varOmega $)}

\begin{table}[!h]
    \centering
    \renewcommand{\arraystretch}{1.5} % 表格行高倍数
    \begin{tabular}{|m{1.3cm}<{\centering}|m{1.3cm}<{\centering}|m{1.3cm}<{\centering}|m{1.3cm}<{\centering}|m{1.3cm}<{\centering}|m{1.3cm}<{\centering}|}
        \hline
        $N$ & $R_s (\varOmega)$ & $R_x (\varOmega)$ & $\Delta R_s (\varOmega)$ & $\Delta n$（格） & $S$（格） \\
        \hline
        0.01 & & & & & \\
        \hline
        0.1 & & & & & \\
        \hline
        1 & & & & & \\
        \hline
        10 & & & & & \\
        \hline
        100 & & & & & \\
        \hline
    \end{tabular}
\end{table}

\begin{tikzpicture}[remember picture,overlay]
    \node[anchor=south east,inner sep=100pt] at (current page.south east) {
        \renewcommand{\arraystretch}{1.5} % 表格行高倍数
        \setlength{\tabcolsep}{18pt}    
    \begin{tabular}{|c|c|}
        \hline
        \LARGE  教师 & \LARGE  姓名 \\
        \hline
        \LARGE \kaishu 签字 &  \\
        \hline
        \end{tabular}
    };
\end{tikzpicture}

\newpage

\section{实验结论及现象分析}


\textbf{对比不同的N值下，惠斯通电桥灵敏度变化，并分析其他可能影响惠斯通电桥灵敏度参量}

由实验数据可知，当N值等于1时，电桥的灵敏度最高；N值越远离1，电桥的灵敏度越低。

其他会影响电桥灵敏度的因素包括：
\begin{enumerate}
    \item \textbf{检流计的灵敏度}：检流计的灵敏度越高，电桥的灵敏度也越高。
    \item \textbf{接线电阻}：接线产生的额外电阻会对电桥的灵敏度产生影响，特别是在高精度测量中。
    \item \textbf{环境温度}：温度变化会导致电阻值发生变化，从而影响电桥的平衡状态和灵敏度。
    \item \textbf{电源电压}：适当提高电源电压可以增大电桥灵敏度。
\end{enumerate}

\section{讨论问题}

\subsection{电桥测电阻为什么不能测量小于$1\Omega$的电阻？}

主要有以下几点原因：

\begin{enumerate}
    \item \textbf{接触电阻和导线电阻的影响}：在测量较小的电阻时，导线本身的电阻以及电桥各接点之间的接触电阻会显著影响测量结果。由于这些电阻通常与待测小电阻相当甚至更大，它们会导致误差增大，难以获得准确的测量结果。
    \item \textbf{检流计灵敏度的限制}：惠斯通电桥的精度依赖于检流计的灵敏度。当电阻非常小时，电桥接近平衡时检流计上的信号太弱，导致检流计上的电流变化幅度非常小，零点不容易准确判断。此时检流计可能无法检测到足够明显的偏差，进而无法精准确定电桥的平衡点，从而降低了测量精度。
    \item \textbf{热效应的影响}：当测量小电阻时，流过电桥的电流较大，可能会导致电路发热，从而使电阻值发生变化，影响测量结果。这种由热效应引起的误差在小电阻测量中尤为突出。
\end{enumerate}

\subsection{用什么方法保护电流计，不至于因电流过大而损坏？}

\begin{enumerate}
    \item \textbf{逐渐调高电流计灵敏度}：在实验过程中，可以先将电流计的灵敏度调至较低水平，然后逐渐增大灵敏度，直到电流计的偏转达到合适的范围。
    \item \textbf{选择合适的电压值}：在实验设计时，选择较小的电压值，以减小对电流计的负载，进一步保护电流计。
    \item \textbf{点按电流计开关}：在接通电流计时通过点按开关的方式，防止大电流持续经过电流计使其损坏。
\end{enumerate}

\subsection{当电桥平衡后，若互换电源和检流计位置，电桥是否仍然平衡？并证明。}

电桥仍然平衡。证明如下：

设电桥的四个臂分别为$R_1, R_2, R_3, R_4$，电桥平衡时，电流计的偏转为0，即：

\begin{equation}
    \frac{R_1}{R_2} = \frac{R_3}{R_4}
\end{equation}

当电源和检流计位置互换后，电桥的平衡条件变为：

\begin{equation}
    \frac{R_1}{R_3} = \frac{R_2}{R_4}
\end{equation}

由于（1）式和（2）式等价，所以电桥在互换电源和检流计位置后仍然保持平衡。

\end{document}